\section{Materials and Methods}
\newcommand{\boxheight}{3.85cm}
\begin{figure}[!h]
	\centering
	\resizebox{\textwidth}{!}{%
		\begin{tikzpicture}[
			font=\sffamily\scriptsize,
			>=Stealth,
			node distance=0.45cm,
			basebox/.style={
				thick,
				rounded corners=4pt,
				inner sep=6pt,
				anchor=center,
				align=center
			},
			phase1/.style={basebox, draw=blue!75!black, fill=blue!4},
			phase2/.style={basebox, draw=violet!75!black, fill=violet!4},
			phase3/.style={basebox, draw=teal!80!black, fill=teal!4},
			phase4/.style={basebox, draw=orange!85!black, fill=orange!4},
			phase5/.style={basebox, draw=green!60!black, fill=green!5},
			arrowline/.style={->, ultra thick, draw=gray!65}
			]
			
			% FASE 1: ML BENCHMARKING (5x5 Repeated NCV)
			\node[phase1] (box1) {
				\begin{minipage}[t][\boxheight][t]{3.15cm}
					\centering
					\textbf{\footnotesize Phase 1: ML Modeling} \par\vspace{2pt}\hrule\vspace{7pt}
					\textbf{Training Set:} \\$N = \NCompounds$ TgDHFR \par\vspace{4pt}
					\textbf{Validation:} \\$5 \times 5$ Repeated NCV \par\vspace{4pt}
					\textbf{Descriptors:} \\ECFP4, 2D, 2D/3D \par\vspace{4pt}
					\textbf{Ensemble:} \\Stacking + RidgeCV \par\vspace{4pt}
					\textbf{Ablation:} \\PI + Gaussian Noise
				\end{minipage}
			};
			
			% FASE 2: FDA SCREENING
			\node[phase2, right=of box1] (box2) {
				\begin{minipage}[t][\boxheight][t]{3.15cm}
					\centering
					\textbf{\footnotesize Phase 2: Virtual Screen} \par\vspace{2pt}\hrule\vspace{7pt}
					\textbf{Target Space:} \\FDA Approved \par\vspace{4pt}
					\textbf{Library Size:} \\$N \approx 1,700$ \par\vspace{4pt}
					\textbf{PhysChem:} MW \\$\le 1000$\,Da \par\vspace{4pt}
					\textbf{Applicability:} \\Isolation Forest \par\vspace{4pt}
					\textbf{Prioritized Set:} \\$N = 19$ candidates
				\end{minipage}
			};
			
			% FASE 3: MOLECULAR DOCKING
			\node[phase3, right=of box2] (box3) {
				\begin{minipage}[t][\boxheight][t]{3.15cm}
					\centering
					\textbf{\footnotesize Phase 3: 3D Consensus} \par\vspace{2pt}\hrule\vspace{7pt}
					\textbf{Receptor:} \\TgDHFR (6AOG, Ch.~\RedockChain{}) \par\vspace{4pt}
					\textbf{Engine:} \\AutoDock Vina \par\vspace{4pt}
					\textbf{Exhaustiveness:}\\ \DockExhaustiveness{} \par\vspace{4pt}
					\textbf{Redock Fidelity:} \\RMSD = \RedockRMSD{}\,\AA \par\vspace{4pt}
					\textbf{Consensus:} \\ML--Physics overlap
				\end{minipage}
			};
			
			% FASE 4: ADMET PROFILING
			\node[phase4, right=of box3] (box4) {
				\begin{minipage}[t][\boxheight][t]{3.15cm}
					\centering
					\textbf{\footnotesize Phase 4: ADMET Profiling} \par\vspace{2pt}\hrule\vspace{7pt}
					\textbf{Input:} \\Docking consensus set \par\vspace{4pt}
					\textbf{Druggability:} \\Lipinski Rule of 5 \par\vspace{4pt}
					\textbf{Route Profiling:} \\Caco-2 (Oral vs. IV) \par\vspace{4pt}
					\textbf{Safety Filter:} \\hERG ($P_\mathrm{block} < 0.5$) \par\vspace{4pt}
					\textbf{Framework:} \\TDC Benchmarks
				\end{minipage}
			};
			
			% OUTCOME / FINAL CANDIDATES
			\node[phase5, right=of box4] (box5) {
				\begin{minipage}[t][\boxheight][t]{3.15cm}
					\centering
					\textbf{\footnotesize Validated Candidates} \par\vspace{2pt}\hrule\vspace{5pt}
					\textbf{Pralatrexate} \par\vspace{1pt}
					\textit{Top Consensus Lead} \par
					$p\mathrm{IC}_{50} = \DockBestMLpIC$ \\$\DockBestDockScore$\,kcal/mol \par\vspace{4pt}
					\textbf{Pyrimethamine} \par\vspace{1pt}
					\textit{High-Efficiency Lead} \par
					$p\mathrm{IC}_{50} = 6.42$ \\$-7.34$\,kcal/mol \par\vspace{4pt}
					\hrule\vspace{3pt}
					\textit{\DockNValidated{} Hits Structurally Verified}
				\end{minipage}
			};
			
			% FLECHAS DE CONEXIÓN
			\draw[arrowline] (box1) -- (box2);
			\draw[arrowline] (box2) -- (box3);
			\draw[arrowline] (box3) -- (box4);
			\draw[arrowline] (box4) -- (box5);
			
		\end{tikzpicture}%
	}
	\caption{\textbf{Schematic overview of the multi-tier computational repurposing and validation pipeline.} 
		The hierarchical funnel integrates: (1) 2D ligand-based representation benchmarking under nested cross-validation, (2) applicability domain-constrained screening of FDA-approved compounds, (3) orthogonal physics-based consensus validation via molecular docking (AutoDock Vina), and (4) in silico ADMET profiling (Caco-2 permeability route categorization and hERG cardiotoxicity filtering) to isolate clinically viable TgDHFR candidates.}
	\label{fig:pipeline_flowchart}
	\vspace{-10pt}
\end{figure}

An overview of the multi-tier virtual screening pipeline implemented in this study is illustrated in Figure~\ref{fig:pipeline_flowchart}.

% ==========================================
% 2. SECTION 2.1
% ==========================================

\subsection{Bioactivity Benchmark Dataset and Curation}

This study utilizes a previously published benchmark of \textit{Toxoplasma gondii} dihydrofolate reductase (TgDHFR) enzyme inhibitors \cite{karimova2025}. By adopting the exact baseline dataset compiled by Karimova and Senanayake, we establish a strict methodological control. This framework isolates the predictive capacity of classical machine learning algorithms from computationally heavy statistical interventions. Deliberately omitting synthetic data augmentation during the primary benchmarking phase ensures that performance metrics reflect empirical ground-truth data rather than algorithmic artifacts.

Initial records merged drug-protein interaction parameters and SMILES structures from BindingDB with bioactivity data from ChEMBL \cite{karimova2025}. Automated standardization protocols, including salt stripping and structural validation, ensured chemical integrity \cite{damilola2025}. To mitigate experimental noise and prevent overfitting \cite{cherkasov2014, iwuala2025}, 117 compounds were systematically excluded from the initial 873 entries. This exclusion targeted organometallic complexes---which preclude 3D force-field generation---and molecules with undefined stereocenters that introduce severe conformational ambiguity.

For compounds exhibiting multiple $IC_{50}$ measurements across varying experimental conditions, we applied the Interquartile Range (IQR) method. Duplicate assays with a bioactivity variance exceeding 1.5 log units were discarded as statistical outliers \cite{karimova2025}. The remaining consistent entries were averaged to yield a single representative measurement per ligand. 

This curation protocol yielded a final benchmark of $N = 756$ unique small molecules. The continuous target variable was transformed into $pIC_{50}$ ($-\log_{10}(IC_{50})$) \cite{karimova2025}. This logarithmic conversion standardizes response values and compresses the active range, optimizing the target distribution for robust regression modeling.

To explicitly quantify the performance impact of synthetic data generation reported in the original literature, we implemented an isolated data augmentation protocol exclusively for post-benchmark ablation studies. Following a strict 15\% hold-out partition to prevent data leakage, training features underwent perturbation via additive Gaussian noise. Applying two distinct noise levels ($\sigma = 0.01$ and $\sigma = 0.001$) generated two synthetic variants per original molecule. This procedure effectively triplicated the training space while maintaining absolute isolation of the empirical test set.

% ==========================================
% 2. SECTION 2.2
% ==========================================

\subsection{Representation Spaces: Molecular Fingerprints and Conformational Descriptors}

To systematically evaluate the structure-activity relationship, the open-source cheminformatics library RDKit \cite{landrum2025} generated four distinct molecular representation spaces. The purely structural baseline utilized Morgan fingerprints (ECFP4), computed with a radius of 2 and folded into a fixed-length vector of \MorganNFeaturesTotal\ bits. Extended-connectivity fingerprints (ECFPs) offer rapid calculation and capture an essentially infinite array of distinct molecular features, including stereochemical topology \cite{rogers2010}. Classical 2D fingerprint-based models frequently achieve predictive parity with state-of-the-art 3D structure-based models when relying exclusively on ligand information \cite{gao2020}.

To establish a quantitative physicochemical baseline, we extracted a comprehensive set of \RdkitTwoDNFeaturesTotal\ conventional two-dimensional topological descriptors. Curated 2D molecular representations demonstrate substantial robustness; recent studies confirm they maintain scaffold-independent advantages over 3D alternatives across various property prediction tasks \cite{tiwari2026}. To evaluate the synergistic potential of combining structural and physicochemical data, a third representation space concatenated the \RdkitTwoDNFeaturesTotal\ 2D descriptors with the \MorganNFeaturesTotal-bit Morgan fingerprints.

To assess the predictive value of spatial geometry, a fourth hybrid representation incorporated 3D conformational descriptors. The Experimental-Torsion Knowledge Distance Geometry (ETKDGv3) approach generated the three-dimensional conformers \cite{riniker2015}. Because purely distance-based constraints often distort aromatic rings and $sp^2$ centers, ETKDGv3 integrates experimental torsion-angle preferences derived from small-molecule crystallographic data to optimize conformation generation \cite{riniker2015}. Following embedding, the MMFF94s force field minimized conformer energy, ensuring thermodynamic stability prior to 3D spatial descriptor extraction. Although recent empirical evaluations suggest that 3D conformer descriptors yield limited incremental predictive value over curated 2D representations in standard tabular benchmarks \cite{tiwari2026}, integrating spatial information remains critical for molecular image-based modeling \cite{matsu2023} and geometry-augmented representation learning \cite{yanan2024}.


% ==========================================
% 2. SECTION 2.3
% ==========================================

\subsection{Nested Cross-Validation Protocol and Variance Correction}

Standard cross-validation partitions datasets into $k$ disjoint blocks, systematically rotating them between training and testing phases \cite{calle2025, krstajic2014}. While conventional nested cross-validation (NCV) literature frequently defaults to symmetrical $5 \times 5$ or $10 \times 10$ architectures \cite{calle2025, lewis2023, wainer2021}, methodological frameworks do not mandate identical fold counts for inner and outer loops \cite{krstajic2014}. To guarantee leak-free evaluation, balance computational overhead, and mitigate model selection bias in this low-data regime ($N = \NCompounds$), we implemented a strict $\TotalOuterFolds \times \InnerKFolds$ asymmetric NCV protocol (\OuterKFolds\ outer folds repeated \OuterRepeats\ times). This architecture mathematically decouples model assessment from model selection via two distinct concentric loops.

Data partitioning strategy dictates the validity of small-scale QSAR modeling \cite{kyrali2022}. Although Bemis-Murcko scaffold splitting is widely advocated to test out-of-domain extrapolation and prevent structural leakage \cite{yang2019, aga2025, beatrice2026}, we deliberately selected a random nested cross-validation scheme. In a highly restricted dataset of \NCompounds\ compounds, scaffold splitting frequently generates severely imbalanced folds. This creates artificial ``activity cliffs'' where entire bioactivity ranges vanish from the training pool, preventing the inner optimization loop from converging on generalizable hyperparameters. Because our primary objective is interpolation within the established applicability domain of known TgDHFR inhibitors---rather than extrapolation to entirely novel chemotypes---random splitting ensures the bioactivity distribution remains statistically consistent across all folds. This consistency enables robust, stable hyperparameter tuning.

The outer loop of the NCV framework evaluates generalization. To achieve lower variance in performance estimation compared to extreme approaches like Leave-One-Out (LOO) \cite{krstajic2014}, we evaluated \TotalOuterFolds\ outer splits. However, repeated cross-validation inherently introduces intra-repetition correlation, violating the strict independence assumption required by standard parametric tests. To rigorously account for this overlapping training data across folds, we applied the Nadeau and Bengio variance correction \cite{nadeau1999} for all downstream statistical comparisons. In each iteration, one fold is strictly held out as the test set while the remaining pool passes into the inner loop. This absolute isolation ensures the test data remains hidden during feature and model selection, yielding a truly unbiased out-of-sample prediction estimate \cite{varma2006, wainer2021}.

The inner loop executes architecture selection and model optimization. Within each outer training pool, an embedded \InnerKFolds-fold inner loop evaluated all single, pairwise, and three-way stacking combinations (up to \MaxComboSize\ base estimators). These combinations were drawn from six diverse regressor families: Random Forest (RF), Extra Trees (ET), Light Gradient Boosting Machine (LightGBM), eXtreme Gradient Boosting (XGBoost), Support Vector Regression (SVR with RBF kernel, $C=10$), and $k$-Nearest Neighbors (kNN with cosine distance). For multi-model combinations, a \texttt{StackingRegressor} utilizing \texttt{RidgeCV} as the meta-estimator coordinated meta-learning. By confining model selection, descriptor scaling (\texttt{StandardScaler}), median imputation, and positive-importance feature selection entirely within the inner partitions, we prevent inadvertent information leakage. This strict encapsulation guarantees that reported metrics reflect genuine generalizability rather than selection artifacts \cite{parvandeh2020, lewis2023, calle2025}.

% ==========================================
% 2. SECTION 2.4 
% ==========================================
\subsection{Deep Learning Baseline: Directed Message Passing Neural Networks}

To benchmark the classical machine learning ensembles against advanced representation learning, we implemented a Directed Message Passing Neural Network (D-MPNN) using the Chemprop framework (version 2.0) \cite{graff2026, yang2019}. To guarantee methodological parity, the D-MPNN underwent the exact \TotalOuterFolds-fold cross-validation protocol utilized for the classical baselines. Executing rigorous resampling on deep learning architectures incurs substantial computational overhead; however, this alignment ensures the resulting performance metrics remain directly comparable and free from single-split variance artifacts.

The architecture processes molecular SMILES strings by mapping them into graph representations via RDKit, designating atoms as vertices and bonds as edges \cite{heid2024}. Initial feature vectors capturing atomic identity and topological connectivity were constructed using \texttt{SimpleMoleculeMolGraphFeaturizer()}. The network applied \texttt{BondMessagePassing()} to extract embeddings directly from directed edges rather than vertices. This edge-centric design eliminates spurious feedback loops (totters), reducing noise in the learned molecular representations \cite{yang2019}. A \texttt{MeanAggregation()} function pooled these atomic embeddings into a unified molecular vector.

Prior to training, the continuous $pIC_{50}$ targets required internal normalization. The aggregated embeddings fed into a \texttt{RegressionFFN()} equipped with batch normalization to yield the final predictive output. Model weights were optimized via PyTorch Lightning over a maximum of 100 epochs. Unlike the classical ensembles, which underwent extensive hyperparameter optimization, the D-MPNN utilized default Chemprop hyperparameters (e.g., hidden size, depth, and learning rate). This methodological asymmetry stems from the prohibitive computational cost of nested cross-validation on graph neural networks using CPU infrastructure, coupled with the high risk of memorization when tuning overparameterized models on small datasets \cite{cawley2010, hadas2023}. Consequently, the D-MPNN represents an unoptimized, out-of-the-box deep learning baseline.

The internal loss function optimized during training was Root Mean Square Error (RMSE). Out-of-sample performance was quantified using the coefficient of determination ($R^2$) and Mean Absolute Error (MAE) averaged across all \TotalOuterFolds\ folds. This rigorous statistical protocol yielded a final cross-validated test performance of $R^2 = \ChemPropDMPNNxCVRTwoMean \pm \ChemPropDMPNNxCVRTwoStd$ and an MAE of \ChemPropDMPNNxCVMae, establishing the deep learning baseline for subsequent comparative analysis.

% ==========================================
% 2. SECTION 2.5
% ==========================================
\subsection{Virtual Screening Protocol and Ligand Efficiency Profiling}

An automated virtual screening pipeline, driven by the highest-performing classical ensemble (\WinnerMode, nested CV $R^2 = \WinnerRTwoMean$), systematically prioritized viable lead candidates from an FDA-approved library. Rather than relying solely on raw predicted bioactivity, the evaluation framework integrated strict applicability domain enforcement with multi-objective ligand efficiency profiling. This approach mitigates the risk of prioritizing ``molecularly obese" compounds that exhibit high theoretical potency but possess unfavorable physicochemical properties for downstream development \cite{hopkins2004}.

Initial structural pre-filtering restricted the FDA library to compounds matching the exact elemental composition of the training set (H, C, N, O, F, P, S, Cl, Br, I). This constraint prevents the model from extrapolating into unsupported chemical spaces. A strict physicochemical ceiling (molecular weight $\le 1000$ Da) controlled for collinearity, preventing molecular size from artificially inflating potency predictions. 

To enforce the applicability domain, an \texttt{IsolationForest} algorithm (5\% contamination threshold) trained on the original Morgan fingerprint matrix excluded structural outliers. This unsupervised anomaly detection ensured all advanced candidates resided strictly within the interpolated chemical space of the known bioactivity distribution. The ensemble's predictive accuracy within this domain was internally validated by accurately recovering the $pIC_{50}$ values of known reference inhibitors, such as Pyrimethamine and Trimethoprim. Ultimately, this hierarchical funnel successfully reduced the initial library to 17 prioritized compounds (10 top-ranked novel candidates and 7 established reference controls) for downstream orthogonal validation.

Candidate prioritization subsequently incorporated four distinct efficiency metrics to correct \textit{in vitro} activity for physicochemical load. Ligand Efficiency (LE) quantified the binding energy per heavy atom count (HAC), requiring candidates to demonstrate optimal atomic utilization \cite{hopkins2014}:
\begin{equation}
	\text{LE} = \frac{1.37 \times pIC_{50}}{\text{HAC}}
\end{equation}
To account for heteroatom contributions and total mass, the Binding Efficiency Index (BEI) substituted HAC with molecular weight (MW) \cite{abad2005}:
\begin{equation}
	\text{BEI} = \frac{pIC_{50}}{\text{MW}} \times 1000
\end{equation}
Furthermore, to balance potency against lipophilicity and mitigate the risk of nonspecific binding, Lipophilic Ligand Efficiency (LLE) was computed as the difference between predicted potency and the calculated partition coefficient ($\text{LLE} = pIC_{50} - \text{cLogP}$). Finally, the Surface Efficiency Index (SEI) evaluated potency relative to the topological polar surface area (TPSA), providing a proxy for membrane permeability potential ($\text{SEI} = (pIC_{50} / \text{TPSA}) \times 100$). Because composite indices combining logarithmic potency values with linear physical parameters can introduce size-dependent mathematical artifacts \cite{shultz2013}, these four metrics (LE, BEI, LLE, SEI) operated strictly as secondary multi-objective ranking filters rather than primary objective functions. 

% ==========================================
% 2. SECTION 2.6
% ==========================================

\subsection{Orthogonal Physics-Based Validation via Molecular Docking}

Orthogonal structure-based evaluation via AutoDock Vina v1.2.0 \cite{rajesh2026, trott2009} complemented the ligand-based predictions. This physical validation translates abstract chemical representations into spatial coordinates, mitigating the risk of false positives induced by algorithmic scaffold memorization. By ensuring genuine structural complementarity, the docking protocol acts as a thermodynamic filter. Vina employs an Iterated Local Search optimizer and a hybrid scoring function, calibrated against PDBbind data, to guide flexible ligands into optimal binding conformations \cite{eberhardt2021}.

Three-dimensional receptor coordinates were retrieved from the Protein Data Bank (\textit{Toxoplasma gondii} DHFR, PDB ID: 6AOG, chain \RedockChain) \cite{consortium2019, berman2003}. To accurately simulate the native biological environment, the NADPH cofactor was retained within the catalytic pocket. Following the removal of crystallographic water and local minimization, polar hydrogens and Gasteiger charges were assigned to generate the required PDBQT files \cite{damilola2025}. The Vina search space was centered on the native \RedockLigand ligand coordinates ($x = \DockCenterX, y = \DockCenterY, z = \DockCenterZ$) with grid box dimensions set to $\DockBoxSize \times \DockBoxSize \times \DockBoxSize\text{ \AA}$ and an exhaustiveness parameter of \DockExhaustiveness. This configuration permitted unrestricted exploration of ligand orientations and flexible rotatable bonds, thereby preventing crystal-pose centering bias \cite{trott2009}.

Prior to screening the machine learning candidates, the docking protocol underwent rigorous validation via re-docking of the co-crystallized \RedockLigand ligand into the TgDHFR active site. The protocol successfully reproduced the native crystallographic binding mode, yielding a top-ranked pose with a centroid distance of $\RedockRMSD\text{ \AA}$. This deviation falls well below the universally accepted $\RedockThreshold\text{ \AA}$ RMSD threshold, confirming the structural fidelity of the grid parameters and the reliability of the scoring function for this specific target.

Following validation, Vina evaluated the 3D conformers of the \DockNTotal\ top-ranked FDA candidates alongside the reference controls. Poses were ranked based on the predicted Gibbs free energy of binding ($\Delta G$) in kcal/mol, where lower energies denote stronger thermodynamic stability \cite{feng2023}. To systematically evaluate the consensus between the ligand-based and structure-based approaches, we established strict quantitative thresholds. Compounds were classified as validated consensus hits only if they exhibited both high ML-predicted potency ($pIC_{50} > 6.5$) and strong docking affinity ($\Delta G < -8.0\text{ kcal/mol}$). Compounds meeting the ML threshold but failing the docking criteria were flagged as potential ML false positives. Finally, to statistically quantify the monotonic agreement between the two orthogonal methods without assuming linearity, a Spearman rank correlation analysis was conducted between the continuous $pIC_{50}$ predictions and the Vina docking scores ($\rho = \DockSpearmanRho$, $p = \DockSpearmanP$).

% ==========================================
% 2. SECTION 2.7
% ==========================================

\subsection{In Silico Pharmacokinetic (ADMET) Profiling}

Following consensus validation, a comprehensive predictive pharmacokinetic (ADMET) assessment evaluated the surviving candidates to reduce late-stage attrition risks. Candidates were first evaluated against Lipinski's Rule of Five criteria, strictly filtering out molecules exceeding 5 hydrogen bond donors, 10 hydrogen bond acceptors, or a calculated LogP of 5. To assess critical toxicity and absorption endpoints, custom machine learning models were trained using curated datasets from the Therapeutics Data Commons (TDC). Molecular structures were encoded as 2048-bit Morgan fingerprints (radius 2) and processed through Random Forest architectures comprising 100 estimators. A \texttt{RandomForestClassifier} trained on the TDC hERG dataset evaluated cardiac toxicity risk, advancing only compounds with a low probability of hERG channel blockade ($P < 0.5$). Concurrently, a \texttt{RandomForestRegressor} trained on the Caco2-Wang dataset predicted intestinal absorption, requiring candidates to demonstrate moderate-to-high permeability ($\log P_{app} > -5.15\text{ cm/s}$). This final ADMET profiling ensured that the prioritized consensus hits represent not only highly efficient binders but also structurally validated, orally bioavailable starting points for therapeutic repurposing.

% ==========================================
% 2. SECTION 2.7
% ==========================================

\subsection{Computational Environment and Reproducibility}

To guarantee algorithmic transparency and absolute computational reproducibility \cite{rojas2026}, all data curation, model training, and evaluation pipelines were programmatically implemented in Python (version 3.14.7). The computational environment relied on a standardized core software stack:

\begin{itemize}
	\item \textbf{Cheminformatics and Featurization:} RDKit (version 2025.9.3) \cite{feng2023} executed all molecular parsing, structural standardization, and descriptor generation (2D, 3D, and ECFP4).
	\item \textbf{Classical Machine Learning:} The nested cross-validation framework \cite{krstajic2014, calle2025}, Bayesian hyperparameter optimization \cite{bauman2014}, and baseline algorithms leveraged \texttt{scikit-learn} (version 1.8.0) \cite{brva2026}, \texttt{xgboost} (version 3.1.2) \cite{wainer2021}, \texttt{lightgbm} (version 4.6.0) \cite{rajesh2026}, and \texttt{optuna} \cite{tyler2025}. Data manipulation and matrix operations were managed via \texttt{pandas} (version 2.3.3) and \texttt{numpy} (version 2.3.5) \cite{eberhardt2021}.
	\item \textbf{Deep Learning Infrastructure:} The D-MPNN graph models utilized \texttt{chemprop} (version 2.0) \cite{graff2026, heid2024} and \texttt{lightning}, with underlying tensor operations supported by PyTorch (\texttt{torch} $\geq$ 2.0.0) \cite{thomas2024}. 
	\item \textbf{Hardware Specifications:} All model training, hyperparameter tuning, and evaluation protocols ran on an Ubuntu 26.04 LTS workstation equipped with an Intel(R) Xeon(R) CPU E5-2695. The deliberate use of CPU-only acceleration explicitly benchmarks the computational accessibility and hardware demands of both classical and deep learning architectures in low-resource settings.
\end{itemize}